\begin{abstract}
Energy efficiency is a paramount concern in Internet of Things (IoT) sensor networks deployed for continuous air quality monitoring. The high sampling rates, multi-channel data acquisition, and persistent wireless transmission required for reliable pollutant tracking impose substantial energy demands on battery- and mains-powered edge devices. This paper proposes a Variance-Aware Adaptive Bit-Width Quantisation (VABQ) framework that jointly optimises sensor-level analogue-to-digital converter (ADC) resolution and inference-level post-training quantisation (PTQ) to minimise energy consumption while preserving monitoring fidelity. The framework exploits the inherent temporal structure of indoor air quality signals---characterised by long stable periods interspersed with activity-driven transient events---to dynamically allocate bit-widths between 4, 8, and 16 bits at the sensor stage, and applies INT8 PTQ to a lightweight one-dimensional convolutional neural network (1D-CNN) for pollutant classification at the inference stage. Experiments are conducted on the DALTON indoor air quality dataset, which spans 28 deployment sites over three months using custom ESP32-based sensor nodes measuring PM$_{2.5}$, PM$_{10}$, CO$_2$, VOC, CO, NO$_2$, temperature, and relative humidity at 1~Hz \cite{karmakar2024exploring}. Results demonstrate that VABQ achieves a 61.3\% reduction in sensor-level ADC energy consumption and a 74.8\% reduction in inference energy relative to full-precision baselines, while maintaining pollutant classification accuracy above 96.1\% and signal reconstruction root-mean-square error (RMSE) below 2.4\% of full-scale range. A comparative analysis against uniform static quantisation and prior adaptive sampling methods confirms the superiority of the proposed dual-layer scheme. The principal novelty lies in the unified optimisation of sensing and inference quantisation through a variance-entropy criterion, bridging two traditionally separate compression domains within a single, mathematically grounded pipeline suitable for resource-constrained edge deployment.
\end{abstract}

\noindent\textbf{Keywords:} sensor quantisation; energy efficiency; air quality monitoring; adaptive bit-width; post-training quantisation; edge AI; Internet of Things; DALTON

\newpage
